\section{Discussion and Conclusions}
\label{sec:Discussion}

The investigation into the Chiral Dirac Equation (CDE) establishes a consistent
algebraic framework for introducing chiral degrees of freedom into relativistic
wave equations. Our analysis confirms that the introduction of the chiral angle
$\alpha$ preserves the fundamental structure of the quantum theory. Explicit
derivation demonstrates that the Hamiltonian $H = \gamma^{0k} p_k + \gamma^0 m
e^{-i \alpha \gamma^5}$ remains strictly Hermitian, ensuring real energy
eigenvalues and unitary time evolution. Consequently, the CDE constitutes a
viable extension of the standard Dirac formalism, physically admissible for the
description of fermionic fields.

\subsection{Chiral Symmetry and Neutrino Phenomenology}

The application of the CDE to neutrino physics does not resolves the vanishing mass
problem inherent in pure left-handed theories. By incorporating a Majorana mass
term, we derived a Chiral Seesaw Mechanism where the effective Dirac mass is
modulated by the chiral phase, $M_D = m e^{-i\alpha}$. This leads to two
distinct physical consequences. First, the active neutrino mass is naturally
suppressed by the heavy sterile scale $M_R$ according to $m'_{\nu} \approx m^2
M_R^{-1} e^{-2i\alpha}$. Second, and perhaps more significantly, the chiral
angle $\alpha$ manifests as a macroscopic observable: it determines the phase of
the light neutrino mass eigenstate. If $\alpha \neq k\pi/2$, this mechanism
provides a geometric origin for leptonic CP violation within the PMNS matrix,
linking the internal symmetries of the Poincar\'{e} group representation
directly to experimental observables beyond the Standard Model.

\subsection{Geometric Scaling of Higher-Spin Masses}

The extension of this formalism to bosonic sectors suggests a unified geometric
scaling law for mass generation. For vector bosons ($s=1$), the enforcement of
chiral symmetry on the constituent spinors of the Bargmann-Wigner multispinor
results in an effective mass $M_{\text{eff}} \approx m \cos^2\left(\alpha\right)$.
Our preliminary analysis of the Rarita-Schwinger field ($s=3/2$) indicates a
continuation of this trend, yielding a cubic scaling $M_{\text{eff}} \approx m
\cos^3\left(\alpha\right)$. This supports the hypothesis that the effective mass
of a particle with spin $s$ scales as $M_{\text{eff}}^{(s)} \approx m
\cos^{2s}\left(\alpha\right)$, identifying mass modulation as a function of the
composite spin structure. However, as noted, the consistency of these
higher-spin equations requires rigorous verification against the Velo-Zwanziger
constraints to preclude acausal propagation modes.

\subsection{Algebraic Generalization via Idempotents}

The construction of wave equations for arbitrary spin $s$ via component-based
tensor calculus becomes intractable for $s > 1$. A more robust approach relies
on the method of Orthogonal Idempotents proposed by Watson and Musielak. Rather
than constructing Lagrangians term-by-term, this formalism utilizes generalized
projection operators $\hat{P}$ to define physical subspaces. For a composite
system of $N$ spinors, the operator is the tensor product of individual helicity
projectors
\begin{equation}
	\hat{P}_{s_1 \dots s_N}\left(\hat{n}\right) = \bigotimes_{i=1}^N
	\hat{P}_{s_i}\left(\hat{n}\right).
\end{equation}
This method is algebraically superior as it handles chiral constraints
intrinsically. By parametrizing the chiral basis of each projector, the
effective mass $M_{\text{eff}}$ emerges not as an ad-hoc parameter, but as an
eigenvalue of the operator acting on the vacuum, $\hat{P}_{\text{total}} \Psi =
M_{\text{eff}} \Psi$. This confirms that chiral degrees of freedom are embedded
in the representation structure of the Poincar\'{e} group. Future work will
utilize this idempotent formalism to classify the causal structure of
the derived higher-spin fields.
